\documentclass[11pt]{article}

\newcommand{\cnum}{Math 132H}
\newcommand{\ced}{Winter 2026}
\newcommand{\ctitle}[4]{\title{\vspace{-0.5in}\cnum, \ced\\Week #1: #2}\author{\vspace{-0.35in}\\#3, #4}}
\usepackage{enumitem}
\usepackage[usenames,dvipsnames,svgnames,table,hyperref]{xcolor}
\usepackage{amsmath, amsfonts}
\usepackage{graphicx} % Required for inserting images
\usepackage{float}
\usepackage{listings}
\usepackage{xcolor}
\usepackage{hyperref}
\usepackage{comment}

\renewcommand*{\theenumi}{\alph{enumi}}
\renewcommand*\labelenumi{(\theenumi)}
\renewcommand*{\theenumii}{\roman{enumii}}
\renewcommand*\labelenumii{\theenumii.}

\definecolor{codegreen}{rgb}{0,0.6,0}
\definecolor{codegray}{rgb}{0.5,0.5,0.5}
\definecolor{codepurple}{rgb}{0.58,0,0.82}
\definecolor{backcolour}{rgb}{0.95,0.95,0.92}
\definecolor{header}{HTML}{205459}
\definecolor{solution}{HTML}{204d52}

\newcommand{\solution}[1]{{{\textcolor{header}{Solution:} \textcolor{solution}{#1}}}}
\setlist[enumerate]{label=(\alph*)}

\lstdefinestyle{mystyle}{
    backgroundcolor=\color{backcolour},
    commentstyle=\color{codegreen},
    keywordstyle=\color{magenta},
    numberstyle=\tiny\color{codegray},
    stringstyle=\color{codepurple},
    basicstyle=\ttfamily\footnotesize,
    breakatwhitespace=false,
    breaklines=true,
    captionpos=b,
    keepspaces=true,
    numbers=left,
    numbersep=5pt,
    showspaces=false,
    showstringspaces=false,
    showtabs=false,
    tabsize=2
}


\begin{document}
\ctitle{\#4}{}{spongebob squarepants}{krusty krab}
\date{}
\maketitle
\section{monday}
Let $U, V$ domains in  $ \mathbb{C}$ or $ \mathbb{C}^{*}$ 
$f : U \rightarrow V$ is a conformal map
if 
\begin{enumerate}
    \item $f$ is holomorphic,
    \item $f$ is surjective
    \item $f$ is one-one
\end{enumerate}
$f^{-1} : V \rightarrow U$ is also conformal (theorem 9.4)\\
Example 1:
Let $U = \mathbb{H} = \{z \in \mathbb{C} : z = x +iy, y > 0\}$ upper half plane.
$V = \mathbb{D} = \{w : |w| < 1\}$ unit disc.
\[
    f: \mathbb{H} \to \mathbb{D} \qquad f(z) = \frac{z-i}{z+i}
\] 
we have $f( \mathbb{H}) \subset \mathbb{D}$ $f$ is one-one and onto from $ \mathbb{C}^{*} \rightarrow \mathbb{C}^{*}$.
and $w = \frac{z-i}{z+i}$ has solutions in $ \mathbb{H}$ if  $|w| < 1$
 \[
     (z+i)w = z-i \qquad z(w-1) = -i(w+1) \implies z = -i \frac{w+1}{w-1}
\] 
and so $z \in \mathbb{H}$ by some argument.
Another example Let
$U = \{1 < |z| < 2\}$ and  $V = \{1 < |w| < 4\}$ and  $f(z) = z^{2}$ $f : U \rightarrow V$ onto but $f$ is  $2-1$ not conformal.
\\Let
$U  = \{1 < |z| < 2\} \cap \{\Im(z) > 0\}$, $f(z) = z^2$ then $f$ is conformal on its image.\\
Example 4.4 Jukovski Map
$z \in \mathbb{D} \rightarrow \frac{1}{2}(z + \frac{1}{z}) = F(z)$ 
$F$ is conformal from $\mathbb{D} \rightarrow f(\mathbb{D})$ and $f(\mathbb{D}) = f( \mathbb{C}^{*} \setminus C\}$ where $C$ is the unit circle
If  $z $ in  $C$ then the image under  $F$ is in the interval $[-1,1]$ so the image of  $ \mathbb{D}$ under $F$ is  $ \mathbb{C}^{*} \setminus [-1,1]$

\section{linear fractional transformations}
Suppose $T : \mathbb{C}^{*} \rightarrow \mathbb{C}^{*}$ conformal (i.e. meromorphic 1-1 onto) so
$Tz = \frac{az+b}{cz + d}$, $a,b,c,d \ne C$ and  $ad - bc \ne 0$
\[
T(\infty) = 
\begin{cases}
    \frac{a}{c} & c \ne 0\\
    \infty & c = 0
\end{cases}
\] 
\[
    L = \{T : Tv \text{ are $LFT$}\}
\] 
\[
    \mathbb{C}^{2} = \{ \begin{pmatrix} z_1 \\ z_2 \end{pmatrix} , z_1,z_2 \in \mathbb{C}\}
\] 
vector space on $ \mathbb{C}$
\[
    \mathbb{C} P^{*} = \{ 1 \text{ (complex) dimensional subspaces of $ \mathbb{C}^{2}$} \}
\] 
let $z_1,z_2 \in \mathbb{C}$ not both $0$ and $[z_1,z_2] = \{ \begin{pmatrix} \lambda z_1 \\ \lambda z_2 \end{pmatrix} : \lambda \in \mathbb{C}\}$
$M : \mathbb{C}P^{*} \rightarrow \mathbb{C}^{*}$ 
\[
= 
\begin{cases}
    \frac{z_1}{z_2}  & z_2 \ne 0\\
    \infty & z_2 = 0
\end{cases}
\] 
$M$ is one-one onto.
Let  $P: \mathbb{C}^2 \rightarrow \mathbb{C}^2$ be a non singuler invertible linear transformation.
so there exists
\[
    \begin{pmatrix} a & b \\ c & d \end{pmatrix} \qquad det \begin{pmatrix}  a & b \\ c & d  \end{pmatrix} \ne 0, a,b,c,d \in \mathbb{C}
\] 
such that
\[
T \begin{pmatrix} z_1 \\ z_2 \end{pmatrix}  = \begin{pmatrix} az_1 + bz_2 \\ cz_1 + d z_2 \end{pmatrix} 
\] 
consider
\[
    T([z_1,z_2]) = [w_1,w_2] \quad \begin{pmatrix} w_1 \\w_2 \end{pmatrix}  = \begin{pmatrix} az_1  + bz_2 \\ cz_1 + dz_2 \end{pmatrix} 
\] 
Let $GL(2, \mathbb{C}) / ( \mathbb{C} \setminus \{0\}) I$ 
$A \sim B \iff AB^{-1} \in ( \mathbb{C} \setminus \{0\}) I$
let 
\[
    \begin{pmatrix} a & b \\c & d \end{pmatrix} \rightarrow^{S} \frac{ az + b}{ cz + d}   \in L
\] 
$S$ is a group isomorphism of   $GL(2, \mathbb{C}) / ( \mathbb{C} \setminus \{0\})I$ into $L$ with  $T,T' \rightarrow T \circ T'$

\section{Tuesday}
\[
    L = \{ Tz = \frac{az + b}{cz + d}, z \in \mathbb{C}^{*}, a,b,c,d \in \mathbb{C}\} = \{f : \mathbb{C}^{*} \rightarrow \mathbb{C}^{*} \text{conformal meromorphic 1-1 onto}\}
\] 
\[
    A = \begin{pmatrix} a & b \\ c& d \end{pmatrix} , B = \begin{pmatrix}  a' & b' \\ c' & d' \end{pmatrix} a'd' - b'c' \ne 0
\] 
\[
T_A = \frac{az + b}{cz + d}, T_B = \frac{a'z + b'}{c'z + d'}
\] 
then
\[
    T_{AB} = T_A \circ T_B
\] 
Corollary:
\[
    T_{A^{-}} = (T_A)^{-1}
\] 
\[
    T_A = I \iff A = \begin{pmatrix} \lambda & 0 \\ 0 & \lambda \end{pmatrix} , \lambda \ne 0
\] 
Corollary:
\[
    T \in L \iff T : \mathbb{C}^{*} \rightarrow \mathbb{C}^{*} \text{ 1-1 onto mermomorphic}
\] 
\[
T'(z) = \frac{ad - bc}{(cz + d)^2} \quad z \ne \infty, Tz \ne \infty
\] 
\[
T(\infty) = \frac{a}{c}
\] 
\[
    T \ne I \implies \text{ T has 1 or 2 fixed points in $ \mathbb{C}^{*}$}
\] 
Proof:
Case 1: $c = 0$, $a \ne d$
 \[
T(\infty) = \infty, \quad Tz = \frac{a}{d}z + \frac{b}{d} \quad ad = 1, a \ne 0, d \ne 0
\] 
then
\[
\frac{a}{d}z + \frac{b}{d} = z \implies z = \frac{-\frac{b}{d}}{1-\frac{a}{d}} = -\frac{b}{a-d} \in \mathbb{C}
\] 
Case 2: $c \ne 0$ then  $T(\infty) =  \frac{a}{c} \ne \infty$
then
\[
Tz = z \iff cz^2 + (d-c)z - b) = 0
\] 
has 1 or 2 roots.\\
4. $L$ is "triply transitive"
\[
    (z_1,z_2,z_3) \in ( \mathbb{C}^{*})^{3} \text{ distinct}
\] 
\[
(w_1,w_2,w_3) \in ( \mathbb{C}^{*})^{3} \text{ distinct}
\] 
then there eixsts unique $T \in L$ such that  $T(z_j) = w_j$  $j = 1,2,3$.
\[
Tz = \frac{z_2-z_3}{z_2-z_1} \frac{z-z_1}{z-z_3}
\] 
if $(w_1,w_2,w_3) = (0,1,\infty)$ since we can assume this without loss of generality( you can compose them).
5. $L$ is the group generated by 
\begin{enumerate}
    \item translations $z \rightarrow z + a$ $\infty \rightarrow \infty$
    \item rotations $z \rightarrow e^{i\alpha}z$ $ \alpha \in \mathbb{R}$ and $\infty \rightarrow \infty$
    \item Dilations $z \to \lambda z$ $\lambda \in \mathbb{R}^{+}, \lambda > 0$ 
    \item Inversion $z \to \frac{1}{z}$
\end{enumerate}
There are 2 cases:
first $c = 0$ then  $\frac{az + b}{d} = |\frac{a}{d}| \frac{(\frac{a}{d})}{|\frac{a}{d}|} z + \frac{b}{d}$
this is a rotation followed by dilation followed by translation.
Case $c \ne 0$
 \[
\frac{az+b}{cz + d} = (\frac{1}{cz + d})(\frac{a}{c}(cz + d) - \frac{ad}{c} + b)
\] 
\[
= \frac{a}{c} + (b - \frac{ad}{c}) \frac{1}{cz + d}
\] 
Define a circle in  $ \mathbb{C}^{*}$ to be any circle in $ \mathbb{C}$ or $\{\infty\} \cup L$ where $L$ is a line in  $ \mathbb{C}$
\section{Theorem 10.2}
$T \in L$  $C$ a circle in $ \mathbb{C}^{*} \implies T(c)$ is a circle  in $ \mathbb{C}^{*}$ \\
Proof:
trivial if $T$ is a rotation, translation, or dilation.
Assume  $Tz = \frac{1}{z}$.
Case 1: $C = L \cup \{\infty\}$ where $L$ is a line $0 \in L$.  $L$ is determined by where it intersects the unit circle at two points
and  $T(L)$ is the line that passes through those points after they are reflected about the origin.
Case 2:  $L$ is a line and  $0 \not\in L$ Let $z^{*} \in L$ with $|z^{*}| = \in\{|z| : z \in L\}$ 
Then $L = \{\infty\} \{z^{*} + \lambda i z^{*}, \lambda \in \mathbb{R}$ then $T(L) = \{0\} \cup \{\frac{1}{z^{*}}\frac{1}{1+ i \lambda}, \lambda \in \mathbb{R}\}$
\[
|\frac{1}{z^{*}} \frac{1}{1+i\lambda} - \frac{1}{2 z^{*}}| = \frac{1}{2|z^{*}|} |\frac{2}{1 + i\lambda} - 1| = \frac{1}{2 |z^{*}|} |\frac{2}{1 _ i\lambda} - \frac{1+ i\lambda}{1 + i\lambda}|
\] 
\[
 =  \frac{1}{2 |z^{*}|} | \frac{1- i\lambda}{ 1 + i \lambda}| = \frac{1}{2 |z^{*}|}
\] 
Case 3 $C = \{|z-z_0| = R\}$, $R > 0$  $z_0 \in \mathbb{C}$
then
\[
    T(c) = \{w : |\frac{1}{w} - z_0|^2 = \mathbb{R}^2\}
\] 
\[
\frac{1}{|w|^2} - 2 \Re( \frac{\overline{ z_0}}{w}) |z_0|^2  + (|z_0|^2- R^2) = 0
\] 
\[
\frac{1}{|z_0|^2 - R^2} - \frac{2}{|z_0|^2 - R^2}\Re(\overline{z_0} \overline{w}) + |w|^2 = 0
\] 
\[
|w-\frac{\overline{z_0}}{|z_0|^2 - R^2}|^2 = \frac{|z_0|^2}{(|z_0|^2 - R^2)^2} - \frac{1}{|z_0|^2 - R^2} > 0
\] 

\section{Wednesday}
Let $\Omega = \mathbb{C}^{*}, \mathbb{C}$ or $ \mathbb{D} = \{z : |z| < 1\}$.
Suppose  $f : \Omega \rightarrow U \subset \mathbb{C}^{*}$ a domain in $ \mathbb{C}^{*}$ and $f$ is conformal,  $f$ meromorphic one-one,  onto
\begin{enumerate}
    \item If $\Omega = \mathbb{C}^{*}$ the n $U = \mathbb{C}^{*}$ then $f = \frac{az + b}{cz  +d}$, $ad -bc \ne 0$ or  $f \in L$.
    \item  $\Omega = \mathbb{C} \implies U \ne \mathbb{C}^{*}$, $p \in  \mathbb{C}^{*} \setminus U$ if $U \ne \mathbb{C}$
        and replace $f$ by $\frac{1}{f(z) - p}$ without loss of generality, $U \subset \mathbb{C}$.
        Now $f : \mathbb{C} \rightarrow U \subset \mathbb{C}$ conformal.
        Consider $\{|z| > R\}$ and  $\{|w| < \frac{1}{R}\}$ $g(w) = f(\frac{1}{w})$ so $g$ is holomorphic on  $\{0 < |w| < \frac{1}{R}\}$ 
        $g$ is holomorphic, open and one-one. The singularity is either removable, essential, or pole. If removable $f$ is bounded and entire therefore constant, not possible.
        If $g$ has an essential singularity it is also impossible since the image in the neighborhood of 0 is dense in the whole plain but the map is open so it cannot be one-one.
        Therefore $g$ has a pole at 0. so
        \[
        g = \frac{1}{w^{N}(a_0 +a_1w + \dots )}
        \] 
        $g$ is $N \to 1$ in a neighborhood of  $0$ therefore $N = 1$ so  $f$ has a pole at infinity simple pole. so $f$ extends  $ \mathbb{C}^{*}$ to $ \mathbb{C}^{*}$ and $\infty \rightarrow \infty$
        so $f(z) = az + b$. Therefore $U  = \mathbb{C}$ 
    \item If $\Omega = \{|z| < 1\}$. Let  $|\alpha| < 1$ and take
        \[
        Tz = e^{i \theta} \frac{z -\alpha}{1-\overline{\alpha}z}
        \] 
        $\theta \in \mathbb{R}$.
        $T\alpha = 0$.
        If $|z| = 1$ then  $|z-\alpha| = |z-\alpha||\overline{z}| = |1- \alpha \overline{ z}|$ and so $ |T(z)| = 1$.
        Let  $M = \{T : Tz = e^{i\theta}\frac{z-\alpha}{1-\overline{\alpha}z}, \theta \in \mathbb{R}, \alpha \in \mathbb{D}\}$
        $M \subset L$
\end{enumerate}

\section{Theorem 10.3}
Schwarz Lemma:
$f : \mathbb{D} \rightarrow \mathbb{D}$ holomorphic and $f(0) = 0$  then  $|f(z)| \le |z|$ and  $|f'(0)| \le 1$.
and equality in either of the above at any point $z_0 \in z$ implies $f(z) = e^{i \theta}z$ for some constant $\theta \in \mathbb{R}$\\
Proof
\[
f(z) = zg(z)
\] 
\[
|f(z)| \le 1 \implies |g(z)| \le \frac{1}{|z|} , \quad |z| < 1
\] 
If $|z| < R < 1$ then $|g(z)| \le \frac{1}{R}$ tends to 1 as $R$ goes to 1,
therefore
 \[
|g(z)| \le 1
\] 
and $g(0) = f'(0)$ and  $|f(z)| \le |z|$.For b equality implies  $|g(z_0)| = 1$ for some $z_0 \in \mathbb{D}$
and so it attains its maximum and so $g$ is constant by maximum principle so  $g(z) = e^{i\theta}$ for some $\theta \in \mathbb{R}$.

\section{Theorem 10.4}
Assume $f : \mathbb{D} \to \mathbb{D}$  $f$ is holomorphic, 1-1, onto- conformal.
Then $f \in M$, i.e. there exists  $\alpha \in D$, $e^{i\theta} \in \partial D$ such that
\[
f(z) = e^{i\theta} \frac{z-\alpha}{1-\overline{\alpha}z}
\] 
Let $\alpha = f(0)$ and set $Tz = \frac{z-\alpha}{1-\overline{\alpha}z}$ and consider $g = T \circ f$ sends  $ 0$ to 0.
Then $|g(z)| \le |z|$ and  $g^{-1} : \mathbb{D} \to \mathbb{D}$ and $g^{-1}(0) = 0$ and
$|g^{-1}(\omega)| \le |\omega|$ by 10.3 so $|g(z)| = |z|$ and  $g(z) = e^{i\theta}z$ 
and so $f = T^{-1}$ and  $ T^{-1} = \frac{z + \alpha}{1 + \overline{\alpha}z}$

\section{Corollary}
Suppose $T: \mathbb{D} \to \mathbb{D}$ is conformal map then
 \begin{enumerate}
     \item $T$ extends to be a 1-1 continuous map from $\overline{ \mathbb{D}} \to \overline{\mathbb{D}}$ 
     \item given $a,b \in \mathbb{D}$ different then there exists  $T \in M$ such that  $T(a) = b$ so  $M$ is transitive
\end{enumerate}

\section{Friday}
Riemann mapping theorem theorem 11.1
\[
    \mathbb{C}^{*}, \mathbb{C}, \mathbb{D} = \{z : |z| < 1\}
\] 
$\Omega \sim \tilde \Omega$ where  $\Omega, \tilde \Omega$ domains in  $ \mathbb{C}$
if there exists $g : \Omega \to \tilde \Omega$ holomorphic 1-1 onto.
If  $\Omega \subset \mathbb{C}^{*}$ is a domain and  $\{ \mathbb{C}^{*} \setminus \Omega$ has at least 2 points
and $\Omega$ is simply connected then there exists $\phi : \Omega \to \mathbb{D}$ conformal onto $ \mathbb{D}$\\
Proof: 
First is a local form of schwarz lemma and the second is properties of normal families. Weierstrass theorem theorems 6.2 and 6.3.
Let
\[
    M = \{ e^{i \theta} \frac{z-\alpha}{1- \overline{\alpha}z} : \theta \in \mathbb{R}, |\alpha| < 1\}
\] 
is the set of conformal mappings from $ \mathbb{D} \to \mathbb{D}$\\
Lemma 11.2:  $f : \mathbb{D} \to \mathbb{D}$ holomorphic. Then
 \[
     |f'(z)|  \le \frac{1- |f(z)|^2}{1 - |z|^2} \quad z \in \mathbb{D}
\] 
and equality holds at some $z \in \mathbb{D}$ if and only if  $f \in M$.\\
Proof: Fix  $z_0 \in \mathbb{D}$ let $S(z) = \frac{z + z_0}{1- \overline{z_0}z}$ then $S(0) = z_0$ and $S \in M$.
Let
 \[
     T(w) = \frac{w - f(z_0)}{1 - \overline{f(z_{0})}w}
\] 
\[
T'(z) = \frac{(1- \overline{\alpha}z) + \overline{\alpha}(z-\alpha)}{(1- \overline{\alpha}z)^2}
\] 
\[
T'(0) = 1 - |\alpha|^2 \quad T'(\alpha)  = \frac{1}{1- |\alpha|^2}
\] 
for general $T \in M$ and in this case swap  $\alpha$ for $f(z_{0})$
let
\[
g(z) = T \circ f \circ S
\] 
so $g(0) = 0$ and  $g : \mathbb{D} \to \mathbb{D}$
 \[
|g'(0)| \le 1
\] 
\[
|g'(0)| = |T'(f(z_0))||f'(z_0)||S'(0)| = |f'(z_0)| \frac{1- |z_0|^2}{1 - |f(z_0)|^2} \le 1
\] 
so we get
\[
|f'(z_0)| \le \frac{1- |f(z_0)|^2}{1 - |z_0|^2}
\] 
equality implies $g = T \circ f \circ S$ is in  $M$ and so  $f$ is in $M$.\\
Normal families
$\{f_n\}_{n=1}^{\infty}$ sequences of holomorphic functions in a domain $\Omega \subset \mathbb{C}$ and $f : \Omega \to \mathbb{C}$ function.
\[
    f_n \to f(z) \quad \text{ uniformly on all compact subsets of $\Omega$}
\] 
implies that $f$ is holomorphic on $\Omega$ and  $f_n^{(k)} \to f^{(k)}$ uniformly on compact subsets of $\Omega$
We clalim  $f_n$ one-one for all  $n $ implies  $f$ is one-one or  $f$ is constant.
Claim 2: If there exists $M < \infty$ such that $|f_n(z)|  \le M$ for all  $n,z$ then  $\{f_n\}$ has a subsequence  $\{f_{n_j}\}$ converging uniformly
on every compact subset $A$ of  $\Omega$.\\
Lemma 11.3: $\Omega, \{f_n\}$ as above then  $f_n$ is 1-1 function,  $f_n \to f$ uniformly in compact  $ K \subset \Omega$ then $f$ is 1-1 or  $f$ is constant.\\
Proof: Assume that $f$ non constant $f(z_0) = f(z_1)$ for $z_0, z_1 \in \Omega$ we can create two discs around $z_1,z_0$ in $\Omega$ and a  $\gamma$ that connects $z_0,z_1$ such that on $\gamma$ $f$ does not equal to  $f(z_0)$ or $f(z_1)$ on $\gamma$
and so we can create a curve surrounding the curve $f(z_0)$ on the map call $\sigma$.
then
\[
\frac{1}{2\pi i} \int_{\sigma}^{} \frac{f'(z)}{f(z) - f(z_0)}dz = 2
\] 
but
\[
\frac{1}{2\pi i} \int_{ \sigma}^{} \frac{f_n'(z)}{f_n(z) - f_n(z_0)}dz = 1
\] 
and the integral converges by uniform convergence on compact domain so this is a contradiction thus $f(z_0) \ne f(z_1)$ \\
Lemma 11.4: 
Let $\{f_n\}$ be a sequence of holomorphic functions on a domain  $\Omega$ and assume there exists a finite number  $M$ such that  $\sup_{\Omega}|f_n(z)| \le M$ for all  $n$.
Then there exists  $(n_j)_{j=1}^{\infty}$ subsequence such that $f_{n_j} \to f$ for some  $f$ converging uniformly on all compact  $K \subset \Omega$
\end{document}
